% Options for packages loaded elsewhere
\PassOptionsToPackage{unicode}{hyperref}
\PassOptionsToPackage{hyphens}{url}
\documentclass[
]{article}
\usepackage{xcolor}
\usepackage{amsmath,amssymb}
\setcounter{secnumdepth}{-\maxdimen} % remove section numbering
\usepackage{iftex}
\ifPDFTeX
  \usepackage[T1]{fontenc}
  \usepackage[utf8]{inputenc}
  \usepackage{textcomp} % provide euro and other symbols
\else % if luatex or xetex
  \usepackage{unicode-math} % this also loads fontspec
  \defaultfontfeatures{Scale=MatchLowercase}
  \defaultfontfeatures[\rmfamily]{Ligatures=TeX,Scale=1}
\fi
\usepackage{lmodern}
\ifPDFTeX\else
  % xetex/luatex font selection
\fi
% Use upquote if available, for straight quotes in verbatim environments
\IfFileExists{upquote.sty}{\usepackage{upquote}}{}
\IfFileExists{microtype.sty}{% use microtype if available
  \usepackage[]{microtype}
  \UseMicrotypeSet[protrusion]{basicmath} % disable protrusion for tt fonts
}{}
\makeatletter
\@ifundefined{KOMAClassName}{% if non-KOMA class
  \IfFileExists{parskip.sty}{%
    \usepackage{parskip}
  }{% else
    \setlength{\parindent}{0pt}
    \setlength{\parskip}{6pt plus 2pt minus 1pt}}
}{% if KOMA class
  \KOMAoptions{parskip=half}}
\makeatother
\usepackage{longtable,booktabs,array}
\usepackage{calc} % for calculating minipage widths
% Correct order of tables after \paragraph or \subparagraph
\usepackage{etoolbox}
\makeatletter
\patchcmd\longtable{\par}{\if@noskipsec\mbox{}\fi\par}{}{}
\makeatother
% Allow footnotes in longtable head/foot
\IfFileExists{footnotehyper.sty}{\usepackage{footnotehyper}}{\usepackage{footnote}}
\makesavenoteenv{longtable}
\setlength{\emergencystretch}{3em} % prevent overfull lines
\providecommand{\tightlist}{%
  \setlength{\itemsep}{0pt}\setlength{\parskip}{0pt}}
\usepackage{bookmark}
\IfFileExists{xurl.sty}{\usepackage{xurl}}{} % add URL line breaks if available
\urlstyle{same}
\hypersetup{
  hidelinks,
  pdfcreator={LaTeX via pandoc}}

\author{}
\date{}

\begin{document}

\section{论文草稿 P0 --- Three-Channel Equivalence
(构造↔直觉↔形式)}\label{ux8bbaux6587ux8349ux7a3f-p0-three-channel-equivalence-ux6784ux9020ux76f4ux89c9ux5f62ux5f0f}

\begin{quote}
版本: 骨架 v0.1 \textbar{} 日期: 2026-08-11 \textbar{}
认识论/数学哲学论文, P1 理论章节 (C5/§1.7) 的独立展开。核心命题见
结构直觉二象性研究.md。
\end{quote}

\begin{center}\rule{0.5\linewidth}{0.5pt}\end{center}

\subsection{1. Introduction}\label{introduction}

\begin{itemize}
\tightlist
\item
  数学基础百年分歧: 直觉主义 (Brouwer: 存在=可构造) / 结构主义
  (Bourbaki: 存在=结构位置) / 形式主义 (Hilbert: 存在=符号系统)。
\item
  现代分歧以类型论 (构造性, Lean/Coq 内核) / 经典数学 (ZFC) / 认知科学
  (System 1/2) 形式存活。
\item
  本文: 在一个神经系统中, 三种立场对应三条真实执行通道,
  对同一输入互可编译地等价 --- 分歧的工程化统一。
\end{itemize}

\subsection{2. Three Channels}\label{three-channels}

\begin{longtable}[]{@{}llll@{}}
\toprule\noalign{}
立场 & 对象 & 通道 & 本项目执行器 \\
\midrule\noalign{}
\endhead
\bottomrule\noalign{}
\endlastfoot
直觉主义 & 构造 & 构造通道 (慢路径) & token\_iterator 定义链展开 \\
结构主义 & 结构位置 & 直觉通道 (快路径) & 编译权重前向 \\
形式主义 & 符号 & 形式通道 (重写) & eval engine / SKI reduce \\
\end{longtable}

\subsection{3. Core Claim}\label{core-claim}

\begin{itemize}
\tightlist
\item
  \textbf{直觉 = 完成了编译的构造 (intuition = compiled construction)}。
\item
  构造提供泛化充分条件 (零构造样本 ⇒ 泛化崩); 定位提供泛化速度 (O(1));
  重写提供可验证性。
\item
  Sem\_构造 = Sem\_直觉 = Sem\_形式: 三通道语义等价 = Computational
  Trinitarianism (Curry-Howard programs=proofs=types) 的神经实现版 ---
  等价发生在同一系统执行通道间, 而非逻辑间。
\item
  反柏拉图主义: 直觉非天赋, 是归纳训练后压缩的构造。
\end{itemize}

\subsection{4. Evidence}\label{evidence}

\begin{itemize}
\tightlist
\item
  零样本曲线 (1-20 样本全 0.0): 无构造 ⇒ 无泛化。
\item
  73 样本学会 super\_log / 2000 位 acc=1.000: 编译后 O(1) 泛化。
\item
  epoch 3 即 1.0: 编译深度单调。
\item
  三通道一致率 {[}待做 EXP-50{]}: 直接证据。
\item
  符号置换 {[}待做 EXP-41{]}: 直觉包裹结构非符号。
\end{itemize}

\subsection{5. Related Work}\label{related-work}

\begin{itemize}
\tightlist
\item
  Brouwer/Bishop/Martin-Löf (构造数学) / Bourbaki (结构主义) / Hilbert
  (形式主义)。
\item
  Curry-Howard / Computational Trinitarianism。
\item
  Kahneman System 1/2。
\item
  LLM 数学能力 vs 数学哲学 (待检索占位)。
\end{itemize}

\subsection{6. Limitations}\label{limitations}

\begin{itemize}
\tightlist
\item
  三通道等价限于已形式化定义的结构 (supported calculus)。
\item
  直觉通道的''结构定位''是统计性, 非保证性 (需验证器)。
\end{itemize}

\subsection{Lean 附录}\label{lean-ux9644ux5f55}

\begin{itemize}
\tightlist
\item
  编译正确性: 直觉通道 = 构造通道的编译 (形式化)。
\item
  重写一致性: 形式通道与定义链展开一致。
\item
  走既有 Lean 纪律 (mathlib 审计/不建独立库/claims ledger)。
\end{itemize}

\begin{center}\rule{0.5\linewidth}{0.5pt}\end{center}

\subsection{结论 (统一框架, 2026-08-11
定稿)}\label{ux7ed3ux8bba-ux7edfux4e00ux6846ux67b6-2026-08-11-ux5b9aux7a3f}

当我们在统一的形式化注意力语法结构下, 最终不可避免地、必然地通过神经网络
清晰、直观、可解释地观察到直觉与构造的边界时, 我们就可以有充足的信心:

\begin{itemize}
\tightlist
\item
  \textbf{用统一的注意力形式语言指导逻辑构造} (形式→构造: 注意力定位
  token 间关系; 形式 ↔ 注意力 attention);
\item
  \textbf{用精确的合成 CoT 逻辑构造获得 flashing 直觉} (构造→直觉:
  合成思维链 编译为快速推理; 构造 ↔ 合成思维链 synthetic CoT);
\item
  \textbf{用蒸馏的直觉上下文穷举注意力形式} (直觉→形式:
  蒸馏快路径搜索组合, 再表达于统一形式; 直觉 ↔ 蒸馏的闪念 distilled
  intuition / flashing).
\end{itemize}

形式、构造、直觉在统一的框架下相互适配,
组成再不分彼此的教学、训练、推理的 反馈迭代管线 --- 这也许是我们通往 ASI
的众多路径中, 相对较快的一条捷径。

\end{document}
